% =============================================================
% บทที่ 3 - วิธีการดำเนินงาน
% โครงงาน: การออกแบบและพัฒนาเกมเดินเรือด้วยระบบสั่งการด้วยเสียงเพื่อส่งเสริมทักษะการสื่อสาร
%          (Design and Development of a Voice-Controlled Ship Navigation Game to Enhance Communication Skills)
% =============================================================

% -----------------------------------------------------------
\section{ขั้นตอนการดำเนินงาน}
% -----------------------------------------------------------

การดำเนินงานวิจัยและการพัฒนาเกมเดินเรือด้วยระบบสั่งการด้วยเสียง แบ่งออกเป็น 7 ขั้นตอนหลัก ดังนี้

\subsection{การวิเคราะห์ปัญหาและศึกษาทฤษฎีที่เกี่ยวข้อง}

\begin{enumerate}[label=(\arabic*)]
    \item ศึกษาค้นคว้าทฤษฎีที่เกี่ยวข้องกับการเรียนรู้โดยใช้เกมเป็นฐาน (DGBL) ทฤษฎีการลดตัวกรองทางอารมณ์ (Affective Filter) และทฤษฎีการเรียนรู้ตามสถานการณ์ (Situated Learning) เพื่อนำมาเป็นกรอบแนวคิดในการออกแบบสภาพแวดล้อมที่ช่วยลดความวิตกกังวลในการพูดภาษาอังกฤษ
    \item ศึกษาคำศัพท์และรูปแบบประโยคมาตรฐานที่ใช้ในการสื่อสารเพื่อการเดินเรือ (Standard Marine Communication Phrases) เช่น การสั่งหันหางเสือ (Port/Starboard) และการสั่งระดับความเร็ว เพื่อนำมาสร้างเป็นคลังคำสั่ง (Command Vocabulary) ภายในเกม
    \item ศึกษาเทคโนโลยีและเครื่องมือสำหรับการพัฒนา ได้แก่ Unity 3D ภาษา C\# รวมถึง API สำหรับการรู้จำเสียง เช่น Whisper AI และการประมวลผลภาษาธรรมชาติ
\end{enumerate}

\subsection{การกำหนดขอบเขตและข้อกำหนดของระบบ}

กำหนดขอบเขตของเกมให้เป็นการจำลองสถานการณ์บังคับเรือหลบหลีกสิ่งกีดขวาง โดยใช้เสียงพูดภาษาอังกฤษผ่านไมโครโฟนเป็นวิธีการหลักในการควบคุมทิศทางและความเร็วของเรือ โดยไม่มีการใช้คีย์บอร์ดหรือเมาส์ในการบังคับทิศทางโดยตรง และกำหนดกลุ่มเป้าหมายสำหรับทดสอบจำนวน 40--50 คน

\subsection{การออกแบบระบบและเกม (Game Design)}

\begin{enumerate}[label=(\arabic*)]
    \item \textbf{การออกแบบเนื้อเรื่องและบทบาท:} ผู้เล่นสวมบทบาทเป็นกัปตันเรือที่ต้องสั่งการลูกเรือ AI (พลต้นหนและพลถือท้าย) เพื่อนำเรือไปถึงจุดหมายอย่างปลอดภัย
    \item \textbf{การออกแบบกลไกหลัก (Core Mechanics):} ผู้เล่นเปล่งเสียงคำสั่ง ระบบแปลงเสียงเป็นข้อความและประมวลผล หากถูกต้องเรือจะเคลื่อนที่ตามหลักฟิสิกส์
    \item \textbf{การออกแบบระบบตอบกลับ (Feedback System):} ออกแบบให้ลูกเรือ AI ทวนคำสั่ง หรือร้องขอให้ผู้เล่นย้ำคำสั่งใหม่เมื่อการสื่อสารติดขัด
    \item \textbf{การออกแบบ UI/UX:} ออกแบบหน้าจอแดชบอร์ด เข็มทิศ เรดาร์ตรวจจับสิ่งกีดขวาง และแถบแสดงสถานะการรับเสียงของไมโครโฟน
\end{enumerate}

\subsection{การพัฒนาเกม (Game Development)}

ดำเนินการพัฒนาเกมโดยแบ่งออกเป็น 4 ระยะ ได้แก่

\begin{enumerate}[label=(\arabic*)]
    \item \textbf{ระยะที่ 1 (Prototype):} พัฒนาระบบพื้นฐาน เชื่อมต่อไมโครโฟนเข้ากับระบบ Speech-to-Text และทดสอบการแปลงเสียงเป็นคำสั่งเพื่อบังคับการเคลื่อนที่ของเรือเบื้องต้น
    \item \textbf{ระยะที่ 2 (Alpha):} พัฒนาสภาพแวดล้อมทางทะเล ระบบสุ่มสร้างสิ่งกีดขวาง และระบบเสียงโต้ตอบของลูกเรือ AI
    \item \textbf{ระยะที่ 3 (Beta):} พัฒนาระบบ UI/UX ให้สมบูรณ์ ปรับจูนระดับความยาก และปรับแต่งความแม่นยำในการรับคำสั่งเสียง
    \item \textbf{ระยะที่ 4 (Release Candidate):} ตรวจสอบและแก้ไขข้อผิดพลาด (Bug Fixing) และจัดเตรียมแพ็กเกจตัวเกมสำหรับนำไปทดสอบกับกลุ่มเป้าหมาย
\end{enumerate}

\subsection{การทดสอบระบบ (Testing)}

ดำเนินการทดสอบระบบเป็น 3 ช่วง ดังนี้

\begin{enumerate}[label=(\arabic*)]
    \item \textbf{ช่วงที่ 1 (Unit Testing):} ทดสอบความแม่นยำและความหน่วง (Latency) ของระบบ Speech-to-Text ในสภาวะแวดล้อมที่มีเสียงรบกวนในระดับต่าง ๆ
    \item \textbf{ช่วงที่ 2 (Integration Testing):} ทดสอบการทำงานร่วมกันระหว่างระบบเสียง ระบบฟิสิกส์เรือ และระบบสิ่งกีดขวาง เพื่อป้องกันข้อผิดพลาดเมื่อผู้เล่นออกคำสั่งที่ไม่คาดคิด
    \item \textbf{ช่วงที่ 3 (User Acceptance Testing: UAT):} ให้กลุ่มเป้าหมายทดลองเล่นเพื่อตรวจสอบความเข้าใจในระบบและค้นหาจุดบกพร่อง
\end{enumerate}

\subsection{การประเมินผล (Evaluation)}

ดำเนินการประเมินผลกับกลุ่มเป้าหมาย 40--50 คน ใน 2 ด้านหลัก:
\begin{itemize}
    \item \textbf{ด้านประสิทธิภาพทางเทคนิค:} วัดค่าความหน่วง (Latency) อัตราความแม่นยำ (Accuracy Rate) และอัตราความผิดพลาด (Error Rate)
    \item \textbf{ด้านประสบการณ์ผู้เล่นและการเรียนรู้:} ประเมินระดับความรู้สึกดื่มด่ำ (Immersion) การลดลงของความวิตกกังวลในการพูดภาษาอังกฤษ และความมั่นใจในการสื่อสาร ผ่านแบบสอบถามและการสัมภาษณ์
\end{itemize}

\subsection{การจัดทำรายงานและสรุปผล}

รวบรวมข้อมูลจากการออกแบบ การพัฒนา ผลการทดสอบ และการประเมินผล มาวิเคราะห์เพื่อจัดทำรูปเล่มรายงานโครงงานฉบับสมบูรณ์

% -----------------------------------------------------------
\section{เครื่องมือที่ใช้ในการดำเนินงาน}
% -----------------------------------------------------------

\subsection{ฮาร์ดแวร์}

รายการฮาร์ดแวร์ที่ใช้ในการพัฒนาและทดสอบโครงงานแสดงดังตารางที่ \ref{tab:hardware}

\begin{table}[H]
\centering
\caption{รายการฮาร์ดแวร์ที่ใช้ในการพัฒนา}
\label{tab:hardware}
\begin{tabular}{|c|l|p{8.5cm}|}
\hline
\textbf{ลำดับ} & \textbf{อุปกรณ์} & \textbf{รายละเอียด/คุณสมบัติ} \\
\hline
1 & เครื่องคอมพิวเตอร์พัฒนา & หน่วยประมวลผลระดับสูง พร้อมการ์ดจอแยกสำหรับรัน Unity 3D และโมเดลปัญญาประดิษฐ์ \\ \hline
2 & ไมโครโฟน (Input Device) & ไมโครโฟนสำหรับรับสัญญาณเสียงพูดภาษาอังกฤษของผู้เล่น \\ \hline
3 & หูฟัง/ลำโพง & สำหรับรับฟังเสียงจำลองสภาพแวดล้อมและเสียงโต้ตอบจากลูกเรือ AI \\
\hline
\end{tabular}
\end{table}

\subsection{ซอฟต์แวร์และบริการ}

รายการซอฟต์แวร์และส่วนเชื่อมต่อโปรแกรมประยุกต์ที่ใช้แสดงดังตารางที่ \ref{tab:software}

\begin{table}[H]
\centering
\caption{รายการซอฟต์แวร์และบริการที่ใช้ในการพัฒนา}
\label{tab:software}
\begin{tabular}{|c|l|l|p{5.5cm}|}
\hline
\textbf{ลำดับ} & \textbf{ซอฟต์แวร์ / บริการ} & \textbf{เวอร์ชัน/ประเภท} & \textbf{วัตถุประสงค์} \\
\hline
1 & Windows 10/11 & OS & ระบบปฏิบัติการหลักสำหรับพัฒนาและรันเกม \\ \hline
2 & Unity 3D & Game Engine & สร้างสภาพแวดล้อม 3D และระบบฟิสิกส์การเดินเรือ \\ \hline
3 & C\# & Programming Language & ภาษาหลักในการควบคุมกลไกภายในเกม \\ \hline
4 & Whisper AI & API / Library & แปลงสัญญาณเสียงพูดภาษาอังกฤษเป็นข้อความคำสั่ง \\ \hline
5 & OpenAI GPT-4o mini & Cloud API & ประมวลผลภาษาธรรมชาติและสร้างบทสนทนาโต้ตอบของ AI\\ \hline
6 & Visual Studio / VS Code & IDE & สภาพแวดล้อมสำหรับการเขียนชุดคำสั่ง \\ \hline
7 & Git / GitHub & Version Control & ระบบควบคุมและจัดการรุ่นของซอร์สโค้ด \\
\hline
\end{tabular}
\end{table}

% -----------------------------------------------------------
\section{แผนการดำเนินงาน}
% -----------------------------------------------------------

แผนการดำเนินงานตลอดระยะเวลาการทำโครงงานในภาคการศึกษาที่ 1 ปีการศึกษา 2569 แสดงดังตารางที่ \ref{tab:gantt}

\begin{table}[H]
\centering
\caption{แผนและระยะเวลาดำเนินการโครงงาน}
\label{tab:gantt}
\begin{tabular}{|c|l|c|c|c|c|c|c|c|c|}
\hline
\textbf{ลำดับ} & \textbf{กิจกรรม} & \multicolumn{2}{c|}{\textbf{ก.ค. 69}} & \multicolumn{2}{c|}{\textbf{ส.ค. 69}} & \multicolumn{2}{c|}{\textbf{ก.ย. 69}} & \multicolumn{2}{c|}{\textbf{ต.ค. 69}} \\ \cline{3-10}
 & & \textbf{1} & \textbf{2} & \textbf{1} & \textbf{2} & \textbf{1} & \textbf{2} & \textbf{1} & \textbf{2} \\ \hline
1 & กำหนดหัวข้อโครงงาน & $\bullet$ & & & & & & & \\ \hline
2 & ศึกษากลุ่มผู้ใช้และงานวิจัย & $\bullet$ & $\bullet$ & & & & & & \\ \hline
3 & กำหนดขอบเขต ข้อจำกัด และเป้าหมาย & & $\bullet$ & & & & & & \\ \hline
4 & การเขียนข้อเสนอโครงงาน & & $\bullet$ & $\bullet$ & & & & & \\ \hline
5 & นำเสนอข้อเสนอโครงงาน & & & $\bullet$ & & & & & \\ \hline
6 & การสร้างและพัฒนาต้นแบบ (Development) & & & $\bullet$ & $\bullet$ & $\bullet$ & & & \\ \hline
7 & การทดสอบ ประเมินผล และปรับปรุง & & & & & $\bullet$ & $\bullet$ & $\bullet$ & \\ \hline
8 & วิเคราะห์และสรุปผลโครงงาน & & & & & & & $\bullet$ & \\ \hline
9 & จัดทำรูปเล่มโครงการ & & & & & & $\bullet$ & $\bullet$ & $\bullet$ \\ \hline
10 & ส่งมอบและนำเสนอโครงงาน & & & & & & & & $\bullet$ \\ \hline
\end{tabular}
\end{table}